\section{Cases}

The propositions earn their place only if they identify real recurrences in subjects that otherwise differ, and the differences are as informative as the recurrences. Five cases follow, and each is analyzed by the same procedure so that the vocabulary does the same work every time. The procedure identifies the transition at issue, says what plays the roles of $C$, $S$, $V$, $P$, $K$ and $t$, names the proposition the case instantiates, and ends with an account of what does not transfer to other domains. The cases are written for this document, and the essays in the corpus that treat the same subjects are cited as sources and are neither excerpted nor summarized, since their terminologies were developed independently and are not assumed to match.

\subsection{Self-distillation}

The setting is a model trained on outputs that it, or a proxy for it, has generated. The proxy-teacher variant is analyzed in an essay in the corpus \cite{contrast}. Three steps are easy to run together and are kept apart here: \emph{selection}, in which some of the generated outputs are kept; \emph{training}, in which the model is refit to what was kept; and \emph{external evaluation}, in which the resulting model is measured against something other than its own outputs. Training on unverified outputs does not mean that no evaluation occurs anywhere in the process, and the absence of an explicit verifier does not by itself establish drift, so the case states one failure mode under stated conditions and shows it without implying that it is general.

The transitions at issue are the inscription of generated outputs into a training record and the commitment of the refit model as the source of the next generation's outputs. The consequence class $C$ is the use of the next generation as a generator or a deployed model. The scope $S$ is the range of items the model is expected to cover. The evidence $V$ is whatever is independent of the model's own outputs, here a known truth. The procedure $P$ is a verifier that rejects incorrect items, or a metric that scores the result, and the commitments $K$ are the earlier generations already relied upon. The time $t$ is the generation index. The model in \texttt{ex3\_selfdistillation.py} is a distribution over thirty correct items, with truth following a $1/\text{rank}$ law, and ten wrong items, refit by counts to one hundred samples per generation for forty generations, averaged over twenty seeds.

Under these conditions, refitting on all samples with no fresh data reduces the proportion of correct items that retain nonzero probability from $1.000$ to $0.127$. Adding a verifier in the loop, which rejects wrong items before refitting, removes the wrong mass entirely and leaves coverage at $0.140$, so the verifier establishes that each retained item is correct and does not establish that the retained items are representative. Adding a tenth as many verified fresh samples from the truth restores coverage to $0.553$. Selecting samples by the model's own confidence reduces coverage to $0.040$ and incidentally removes the wrong mass, since the wrong items start with low probability. A fourth condition computes the same metrics each generation without altering anything and reproduces the first condition's trajectory exactly on every seed, which is the sense in which evaluation that only measures is distinguished from evaluation that changes selection or training. The average wrong mass in the unfiltered condition rises from $0.080$ to $0.111$, but the simulation was not designed to test amplification of error and this figure is not interpreted as evidence of it.

The case instantiates the fifth proposition. Each item in the verified record is individually warranted and the record as a whole is not a faithful account of what it was meant to cover, so the number and correctness of inscriptions underdetermines whether the record can serve its consequence. It also illustrates the difference between a verdict about an item and a verdict about a body of items, which is a difference of scope in the sense of the second proposition.

What does not transfer is considerable. The model is a categorical distribution refit by counts, and nothing about gradient-based training, regularization, architecture or the mixing of real data follows from it. The failure mode holds under the stated conditions and is not claimed for any particular system. Coverage is defined here relative to a known truth, whereas in practice the independent criterion that would make coverage measurable is typically what is unavailable, and that unavailability is itself a limit on external evaluation and not a detail of the toy.

\subsection{Corrective Navigation and the Repair Ledger}

The setting is a navigable record in which corrections are themselves entries, as in the corrective-navigation specification in the corpus \cite{corrective}. A repair of an earlier entry is inscribed like any other, and the question is what gives the repair standing over the entry it repairs, since an inscription is not an evaluation and a later entry is not for that reason a better one.

The transitions at issue are the inscription of a correction and the evaluation that would let it supersede its target. The consequence class $C$ is the replacement of the earlier entry, for readers who reach the region of the record through navigation, by the correction. The scope $S$ is the region of the record that the correction claims to repair. The evidence $V$ is whatever shows both that the target was in error and that the correction fixes the error without introducing another. The procedure $P$ is the process by which corrections are accepted, which may be the same as the process for original entries or stricter. The commitments $K$ include everything already relied upon that was built on the target. The time $t$ matters because a correction accepted early may be revised by evidence available later.

Two propositions apply. The third bears on the acceptance of corrections by a repair procedure. A repair procedure that has accepted many corrections has an acceptance record, and by the third proposition that record cannot show that the procedure repairs, since a procedure that accepts every proposed correction produces the same record as one that accepts only correct ones until a bad correction is tested against something independent. The fifth bears on the ledger as a whole. A record to which corrections are continually appended grows in inscriptions, and whether it grows in warranted content depends on whether each correction is grounded in evidence about the item it replaces or only cites a previous entry that cites another, in the cyclic arrangement the fifth proposition exhibits. A record whose repairs refer only to one another has the appearance of continuous repair and the warrant of none.

What does not transfer is the idea that the same correction procedure suits every domain. A correction to a mathematical derivation is checked against rules of inference and can be settled finally. A correction to an empirical claim is a revision of evidence and is provisional. A correction to an institutional record, such as an erratum to a ruling, is an act of an authority and is valid by virtue of that authority. A repair ledger that records all three in one format has not thereby made them the same kind of thing, and the standing of each is established differently.

\subsection{Economic Authorization without Effective Accountability}

The setting is a decision about production or allocation that has been authorized by a collective procedure of the kind proposed under the heading of economic democracy, and the source in the corpus is a response to an interview on that subject \cite{authz}. The case concerns structure only and takes no position on whether any particular arrangement achieves what it intends. The transition at issue is authorization, and the claim is that the standing of the authorization and the standing of the consequences claimed for the decision are two different things established on different evidence.

The consequence class $C$ is the allocation's effects on those who bear them. The scope $S$ is the population affected, which may be wider than the population that took part in the procedure. The evidence $V$, for the authorization, is the record of how the decision was made and whether the procedure was followed, and for the consequences, it is information about what the allocation actually does. The procedure $P$ for the authorization is the collective decision rule, and the procedure for the consequences is whatever verification is applied to outcomes. The commitments $K$ are prior decisions, contracts and dependencies that the new decision inherits. The time $t$ matters because outcomes arrive after authorization.

The first proposition applies directly. A decision can be licensed by the rules of the collective procedure and lack standing for the consequences it is said to have, and a decision can be defensible in its consequences while the procedure that produced it was irregular. The case is an instance of the occurrence of a licensed transition not entailing its warrant, and of the licensing not being the same as the standing for any consequence that the procedure does not examine. The second proposition adds that standing for a population that participated need not extend to a population that did not, so a decision that has institutional standing within its scope may lack it when its effects fall beyond. The structure of the claim is that verification of outcomes is a different operation from authorization of decisions, and a system that carries out the first well has not thereby carried out the second.

What does not transfer is the normative content of the procedure. Democratic legitimacy is not merely a warrant relation, and a procedure may have a value for the participants that is independent of how good its outcomes turn out to be. The framework says that procedural authorization does not establish outcome standing, and it does not say that procedural authorization is merely instrumental. It also makes no empirical claim about how often decisions made by collective procedures succeed or fail.

\subsection{Reconstructive Memory}

The setting is a store in which retrieving an item regenerates it at the moment of use, and the regenerated item is then recorded or relied upon as though it had been stored and fetched. Experimental work on human recollection has long shown that recalled content is rebuilt from fragments and expectations, and can be altered by the circumstances of recall \cite{bartlett,loftus}. A computational analogue is a memory system that holds an index and summaries and generates the text it returns. The case is written last of the five and checked against the definitions of the second section, because it presents the greatest risk of confusing categories.

The transition at issue is an inscription that carries a label, for instance that the item was retrieved, which misdescribes the transition that took place, which was a generation. This is first a failure of provenance. The record states that one kind of transition occurred when another did, and the first proposition applies, since what occurred and what the record says occurred are different matters. Whether the regenerated item has epistemic standing is a separate question, answered by further evidence. The consequence class $C$ is reliance on the item as a basis for action or testimony. The scope $S$ is the situation in which it is used. The evidence $V$ is corroboration from independent records or sources, and the procedure $P$ depends on the sense: in the epistemic sense a comparison with corroboration, in the computational sense a check of the regenerated text against stored source if there is one, and in the institutional sense the rules of a body that admits testimony. The commitments $K$ are the actions already taken on the item, and the time $t$ is the time of recall, which matters because reconstruction can differ across recalls.

Three assessments are kept apart. The mislabel is a provenance failure and is cured by correcting the label, which does not touch the content. The epistemic standing of the content depends on corroboration and may be high or low, and that reconstruction is compatible with accuracy is not in doubt. A computational check that the regenerated text matches a stored source establishes a match and not the truth of the source. An institutional decision to admit the item as testimony establishes admission and not correctness. A failure in one of these is not a failure in another, and presenting all four as one failure, whether of verification or of authorization, loses what the case teaches.

What does not transfer is any estimate of how reliable reconstructed items are. The literature on human recollection is varied, and it is cited only to show that regenerated content can be inscribed as retrieved. Nothing here implies that reconstructed recollections lack standing in general, or that stored items have it in general: in each direction the standing has to be established, not presumed.

\subsection{Irreversible Commitment in a Branching System}

The setting is distributed version control, in which a history is a graph of commits and branches are named pointers into it. The branching-coordination essay in the corpus treats the setting at length \cite{echo}, and the case here concerns one operation in it. The transition at issue is a rewrite of history that squashes several commits into one, and the commitment at issue is the publication of some of those commits so that collaborators may base work on them.

The consequence class $C$ is the continued use of the shared history by collaborators. The scope $S$ is the set of clones and branches that contain the affected commits. The evidence $V$ is that the final tree is unchanged by the rewrite, and the procedure $P$ accepts a rewrite when that holds and when no commit that is published is removed. The commitments $K$ are the published commits, and the time $t$ is the time of assessment. The program \texttt{ex2\_branching.py} evaluates the same rewrite, with the same evidence, in five states that differ only in which commits are published. The rewrite is accepted when nothing is published or only the base commit is, and is rejected when any removed commit is published. A stateless assessment, which looks at the operation and the tree evidence alone, accepts in all five, and so is mistaken in three. Publishing a commit is a commitment transition, and after it the rewrite that was acceptable is not.

The case instantiates the second and fourth propositions together. By the second, the commitment changes what the rewrite would newly make irreversible, which is why standing attaches to the transition in its state. By the fourth, the rewrite is reversible in the record, since the old objects remain stored and can be restored by their holder, and is not recoverable, since collaborators' copies have diverged and must be repaired by each of them. The standard documentation of the tools advises against rewriting history that others have based work on for exactly this reason \cite{gitrebase}. The first section's remark on de facto commitment also applies, since a commit that collaborators have begun to rely on is published in the relevant sense whether or not anyone designated it so.

What does not transfer is the clarity of the dependency structure. In version control the dependencies are explicit, because commits are identified by content and clones can be enumerated in principle, so what is published is in principle observable. In most other domains who relies on what is not recorded and has to be reconstructed, and much of the difficulty of commitment lies in that reconstruction. The rule against rewriting published history is also a convention with exceptions, such as private branches and coordinated rewrites that all collaborators have agreed to, and its force is relative to what a given community has committed to and not an absolute property of the operation.
